\section{The ordinary scaling construction}\label{sec:scaling}

We record the ordinary scaling construction used in the
lower-dimensional induction of Section~\ref{sec:contractions}. The
construction starts on the given pair and uses only ordinary negative
steps, with a fixed Kähler class as the scaling direction.

Fix a smooth Kähler form \(h\) on \(X\), with class \(H\), large enough
that \(K_X+\Delta+H\) is nef.  Throughout this section
\(\mathbf H\) denotes the fixed b-nef datum obtained by pulling back
this very form.  On subsequent models its trace class is denoted by
\(H_T\), as in Lemma~\ref{lem:class-trace}.  We do not replace the
fixed datum by arbitrary cohomologous current representatives.

\subsection{Constructing the scaling program}

\begin{proposition}\label{prop:scaling-construction}
Let \((X,\Delta)\) satisfy the hypotheses of Theorem~\ref{thm:finite}.
There is an ordinary negative-ray program, finite or infinite,
\[
 (X,\Delta)=(X_0,\Delta_0)\dashrightarrow
 (X_1,\Delta_1)\dashrightarrow\cdots,
\]
in the same global Weil-divisor \(\Q\)-factorial compact Kähler
category.  It stops when the ordinary adjoint is nef or a negative ray
has a Mori fibre contraction.  Put
\[
 D_i=K_{X_i}+\Delta_i,\qquad
 L_i(t)=D_i+tH_i,\qquad H_i=H_{X_i},\qquad \lambda_{-1}=1.
\]
At every stage before stopping there is a parameter
\(0<\lambda_i\leq\lambda_{i-1}\) and a contracted extremal ray
\(R_i\) such that
\[
 D_i\cdot R_i<0,\qquad L_i(\lambda_i)\cdot R_i=0,
 \qquad L_i(\lambda_i)\text{ is nef}.
\]
The generalized pairs \((X_i,\Delta_i+t\mathbf H)\) are gklt for
\(0\leq t\leq\lambda_{i-1}\).  A birational step is crepant for the
data at \(t=\lambda_i\), and discrepancies do not decrease across
that step for any \(0\leq t\leq\lambda_i\).
\end{proposition}

\begin{proof}
On \(X_0\), the nef datum descends, so its addition changes no
discrepancies.  Thus all the initial generalized pairs are gklt.
Suppose the assertions have been established through \(X_i\).
If \(D_i\) is nef, stop.  Otherwise define
\[
 \lambda_i=\min\{t\in[0,\lambda_{i-1}]:L_i(t)\text{ is nef}\}.
\]
The set is nonempty by induction and is closed and convex.  Its minimum
is positive because \(D_i\) is not nef.

We first find a ray at this threshold.  Choose
\(t_k\uparrow\lambda_i\) with \(\lambda_i/2<t_k<\lambda_i\).
The generalized cone theorem gives an extremal ray negative for
\(L_i(t_k)\).  Since \(L_i(\lambda_i)\) is nef, that ray is also
negative for \(L_i(\lambda_i/2)\).  The trace \(H_i\) is big by
Corollary~\ref{cor:big-trace}.  Hence
\(\Delta_i+(\lambda_i/2)H_i\) is big, and the cone theorem gives
only finitely many negative rays for this latter generalized pair
\cite[Theorem~1.3]{HX2026}.  One ray \(R_i\) therefore occurs for
infinitely many \(k\).  On this fixed ray, continuity gives
\(L_i(\lambda_i)\cdot R_i=0\).  Moreover
\(H_i\cdot R_i>0\), so \(D_i\cdot R_i<0\).

Apply Proposition~\ref{prop:ordinary-step} to this ordinary negative
ray, and denote its contraction by \(f_i:X_i\to Z_i\).  A fibre
contraction gives the stopping alternative.  Otherwise make the
ordinary divisorial step or flip.  Lemma~\ref{lem:class-transport}
transports the trace \(H_i\) to \(H_{i+1}\) while preserving its
relation to the fixed b-datum.  By Lemma~\ref{lem:pullback-image}, the
class \(L_i(\lambda_i)\) descends to \(Z_i\).  Its descended class
is nef, and its pullback on the new model is
\(L_{i+1}(\lambda_i)\).  Thus the two threshold classes have equal
pullbacks on a common resolution.  Their difference is represented by
the actual exceptional discrepancy divisor, which vanishes by
negativity as in Lemma~\ref{lem:small-comparison}.  The step is
consequently generalized crepant at the threshold.

For completeness, this also proves the induction on singularities.
On a common resolution, the log discrepancy at each prime \(E\) is
affine in \(t\).  Its difference across the step is zero at
\(t=\lambda_i\) and is the ordinary discrepancy increase at zero.
Consequently
\begin{align}
 &a(E,X_{i+1},\Delta_{i+1}+t\mathbf H)
       -a(E,X_i,\Delta_i+t\mathbf H)\notag\\
 &\qquad=\left(1-\frac{t}{\lambda_i}\right)
     \bigl(a(E,X_{i+1},\Delta_{i+1})-a(E,X_i,\Delta_i)\bigr)
     \geq0\label{eq:scaling-interpolation}
\end{align}
for \(0\leq t\leq\lambda_i\).  The generalized pairs on this
interval remain gklt, and the new upper-threshold class is nef.  The
ordinary steps preserve all the stated properties of the original
category by Proposition~\ref{prop:ordinary-step}.
\end{proof}
